\scanpage{24}
d'où

\[
\partial_n:\underline H^1_Y(G_{mX})\longrightarrow\underline H^2_Y(\unsure{\mu_{n_X}|\mu_{nX}})
\]

et nous poserons

\[
\deleted{\unsure{\xi|\varphi}}\ \underline\gamma^{(n)}_{Y/X}=\partial_n(\underline\beta_{Y/X})\quad\in\Gamma(X,\underline H^2_Y(\mu_{nX}))
\]

en laissant tomber \unsure{le|les} exposant $n$ s'il n'y a pas d'ambiguité possible.

\inserted{La définition de} Cette section \deleted{satisfait} commute à la localisation de Zariski (et $=$ à tt \unsure{changt|changemt} de base $X'\longrightarrow X$ tel que $Y'$ reste une ``int. complète'' i.e. $\varphi'$ reste section régulière) ce qui permet de définir $\gamma^{(n)}_{Y/X}$ en général (sans équation globale). Si \inserted{maintenant} $Y$ est intersection complète de codim $d$, alors localemt $Y$ est de la forme

\[
Y=Y_1\cap\ldots\cap Y_d
\]

et on posera

\[
\underline\gamma_{Y/X}=\underline\gamma^{\deleted{\unsure{(n)}}}_{Y_1/X}\mathbin{.}\ \mathbin{.}\underline\gamma_{Y_d/X}\quad\in\Gamma(Y,\underline H^{2d}_Y(\mu_n^{\otimes d}))
\]

(\uline{cup-produit}), et on prouve que cette section ne dépend pas du choix des $Y_i$, donc se définit globalement.

Propriétés générales :

% archive sheet number: 23
